\section{Conclusion}

We probe the hidden states of two architecturally distinct robot policies and show they encode failure risk and temporal urgency---extending detection-only approaches like SAFE~\citep{safe_paper} to include time-to-failure prediction and correction estimation. We attempted to use these signals for continuous spatial steering, but found that the learned correction direction does not outperform random perturbation. The detection timing is what matters: knowing \emph{when} to intervene is the operative mechanism, not \emph{where} to push.

PULSE reads these signals with a lightweight probe at $6{\times}$ lower latency than MPPI, self-calibrating to difficulty and generalizing zero-shot. The jiggle ablation simplifies the safety-adapter design space and identifies correction fidelity as the clear bottleneck.

\paragraph{Improving the correction head.}
Three directions may address the correction bottleneck:
\begin{enumerate}[leftmargin=1.5em, topsep=2pt, itemsep=1pt]
    \item \textbf{Recovery-conditioned labels.} Current labels use nearest-neighbor trajectory matching ($\Delta\mathbf{p}_t = \mathbf{p}^{\mathrm{succ}}_t - \mathbf{p}^{\mathrm{fail}}_t$). Training on demonstrations that \emph{recovered from near-failure} would provide cleaner directional signal.
    \item \textbf{Multi-step correction planning.} Single-step Cartesian deltas cannot capture multi-step recovery. A sequence model over hidden states could plan corrective trajectories.
    \item \textbf{Task-specific correction heads.} Training on failure modes of specific tasks and embodiments rather than a universal head.
\end{enumerate}

\paragraph{Additional future work.}
{\color{blue}(i) Physical robot validation on SO-101 arms with multiple LeRobot VLAs.}
(ii) Nonlinear cross-model transfer.
(iii) Integration with conformal prediction guarantees for calibrated failure detection.

All code, data, checkpoints, and evaluation scripts will be released.
